% =====================================================================
%  novelty_and_contributions.tex  --  PHASE I, Task 2
%
%  Reframes the paper from an application note ("we apply GMM and LSTM")
%  to a methodology contribution ("we propose a unified framework").
%
%  Contains, in order:
%    1. Reframed abstract               -> replaces the current abstract
%    2. Positioning / novelty statement -> end of Section 1 (Introduction)
%    3. Itemised contribution list      -> immediately after (2)
%    4. Paper organisation paragraph    -> closes Section 1
%
%  Drop-in usage: copy the relevant block into the main document, or
%  % =====================================================================
%  novelty_and_contributions.tex  --  PHASE I, Task 2
%
%  Reframes the paper from an application note ("we apply GMM and LSTM")
%  to a methodology contribution ("we propose a unified framework").
%
%  Contains, in order:
%    1. Reframed abstract               -> replaces the current abstract
%    2. Positioning / novelty statement -> end of Section 1 (Introduction)
%    3. Itemised contribution list      -> immediately after (2)
%    4. Paper organisation paragraph    -> closes Section 1
%
%  Drop-in usage: copy the relevant block into the main document, or
%  % =====================================================================
%  novelty_and_contributions.tex  --  PHASE I, Task 2
%
%  Reframes the paper from an application note ("we apply GMM and LSTM")
%  to a methodology contribution ("we propose a unified framework").
%
%  Contains, in order:
%    1. Reframed abstract               -> replaces the current abstract
%    2. Positioning / novelty statement -> end of Section 1 (Introduction)
%    3. Itemised contribution list      -> immediately after (2)
%    4. Paper organisation paragraph    -> closes Section 1
%
%  Drop-in usage: copy the relevant block into the main document, or
%  % =====================================================================
%  novelty_and_contributions.tex  --  PHASE I, Task 2
%
%  Reframes the paper from an application note ("we apply GMM and LSTM")
%  to a methodology contribution ("we propose a unified framework").
%
%  Contains, in order:
%    1. Reframed abstract               -> replaces the current abstract
%    2. Positioning / novelty statement -> end of Section 1 (Introduction)
%    3. Itemised contribution list      -> immediately after (2)
%    4. Paper organisation paragraph    -> closes Section 1
%
%  Drop-in usage: copy the relevant block into the main document, or
%  \input{novelty_and_contributions} after \section{Introduction}.
%  Bibliography: paper/references.bib
% =====================================================================


% ---------------------------------------------------------------------
% 1. REFRAMED ABSTRACT
% ---------------------------------------------------------------------
% Replaces the current abstract. Structure: context -> gap -> what we
% propose (three named components) -> how it is validated -> headline
% result -> significance. Numeric placeholders marked XX must be filled
% from the final experimental run before submission.
%
% PHASE III WARNING ON THE FINAL SENTENCE OF RESULTS.
% The placeholder "...at a payback period of XX years" cannot be filled
% with a favourable figure for pelletizer-I. The plant already
% de-energises for the long idle gaps, leaving 86.8 h of energised idle
% over 63 days of coverage (~325 kWh, ~39 USD/yr at 0.12 USD/kWh), so a
% 5,000 USD retrofit pays back in ~130 years and a three-year payback
% needs a per-machine cost under ~116 USD. Do NOT quote a payback number
% as a headline. The defensible claims are (a) recoverable energy, and
% (b) decision quality: the proposed policy captures 55-76% of the
% attainable head-room and removes the status quo's 20 constraint-
% violating and 29 loss-making shutdowns. Payback belongs in the
% sensitivity discussion. Revisit after Phase IV establishes whether any
% IMDELD machine has a larger energised-idle fraction.
% ---------------------------------------------------------------------

\begin{abstract}
Industrial machines consume a substantial fraction of their total
electrical energy while energised but not producing, yet eliminating
this standby consumption in an existing facility is obstructed by three
coupled difficulties: per-timestamp operating-state annotations do not
exist and are prohibitively expensive to obtain; the duration of a
future non-productive interval is unknown at the moment a shutdown
decision must be made; and shutting a machine down is only economical
when that duration exceeds a break-even point set by restart energy,
tariff, and production constraints. Existing work addresses these in
isolation---non-intrusive load monitoring infers states but does not
forecast or act, while energy-aware scheduling decides optimally but
assumes the idle duration is given. This paper proposes a unified,
label-free framework that closes the loop from raw electrical
measurements to an economically justified shutdown decision. The
framework comprises three components: (i)~\emph{physics-validated
unsupervised state inference}, in which a Gaussian hidden Markov model
jointly learns the emission distributions and the transition dynamics of
the machine's operating states, and the resulting partition---which
cannot be checked against ground truth---is validated by an ensemble of
independent physical, circuit-level, and operational-schedule tests
rather than by supervised metrics; (ii)~\emph{non-productive duration
forecasting}, in which a sequence-to-sequence LSTM predicts the length
of the upcoming standby interval from the inferred state history; and
(iii)~\emph{optimisation-based shutdown decision support}, in which the
break-even duration is derived as the solution of a constrained
cost-minimisation problem over standby, restart, and production-delay
costs, rather than supplied as a tuning parameter. The framework is
evaluated end-to-end on IMDELD, a public benchmark of 1\,Hz
measurements from eight heavy machines in a full-scale industrial
plant. Temporal modelling reduces the state-sequence flicker rate from
XX\% to XX\%, the forecasting stage attains a standby-duration MAE of
XX, and the decision layer identifies XX\,kWh/year of recoverable
standby energy per machine at a payback period of XX~years. The
framework requires only the voltage and current measurements a standard
retrofit meter already provides, and its outputs map directly onto
ISO~50001 energy performance indicators.
\end{abstract}


% ---------------------------------------------------------------------
% 2. POSITIONING / NOVELTY STATEMENT
%    Place at the end of Section 1, immediately after the gap has been
%    stated. This is the sentence reviewers will quote back.
% ---------------------------------------------------------------------

% --- BEFORE (do not use; retained here only to document the change) ---
% "We apply a Gaussian Mixture Model for state identification and an
%  LSTM for forecasting."
%  -> Reads as an application of two off-the-shelf models. Invites the
%     reviewer question "what is new here?" with no defensible answer.
% ----------------------------------------------------------------------

\medskip
\noindent
To close this gap we propose a \emph{unified data-driven framework for
industrial standby energy elimination}. Rather than treating state
inference, forecasting, and control as separate problems, the framework
treats them as one pipeline whose stages are designed against each
other's requirements: the inference stage is constrained to produce a
temporally coherent sequence because duration forecasting is otherwise
ill-posed, and the forecasting stage is constrained to predict duration
rather than class because the decision stage consumes a duration. The
framework consists of three components:

\begin{enumerate}
  \item \textbf{Physics-validated unsupervised state inference.} A
  Gaussian hidden Markov model \cite{rabiner1989tutorial} with a sticky
  self-transition prior \cite{fox2011sticky} jointly estimates the
  emission distributions and the transition dynamics of the machine's
  operating states directly from electrical measurements, requiring no
  annotation. Because no ground-truth labels exist against which the
  resulting partition could be scored, we replace supervised evaluation
  with a validation ensemble drawn from sources external to the model:
  induction-motor power-factor and no-load current physics, conservation
  of power across the plant's circuit hierarchy, dwell-time
  plausibility, cross-machine replication between structurally identical
  units, and agreement with the factory's known shift schedule.

  \item \textbf{Non-productive duration forecasting.} A
  sequence-to-sequence LSTM \cite{hochreiter1997lstm,sutskever2014seq2seq}
  maps the recent inferred-state history to the duration of the upcoming
  standby interval. Predicting duration, rather than the next state
  label, is what makes the output actionable: the break-even criterion
  is a statement about how long a machine will remain idle, not about
  what it is doing now.

  \item \textbf{Optimisation-based shutdown decision support.} The
  shutdown decision is posed as minimisation of total cost---standby
  energy, restart energy and labour, and production-delay penalty---
  subject to thermal cool-down, restart-frequency, and schedule
  constraints, in the spirit of \cite{mouzon2007operational} but with
  the idle duration supplied by component~(ii) instead of by a known
  schedule. The break-even duration is therefore a \emph{derived}
  quantity of the optimisation rather than an operator-supplied
  threshold, and its sensitivity to tariff and restart energy is
  reported explicitly.
\end{enumerate}

\noindent
The framework is validated end-to-end on the IMDELD benchmark
\cite{imdeld2018dataset,martins2018imdeld}, comprising 1\,Hz electrical
measurements from eight heavy machines in a full-scale Brazilian
poultry-feed pellet plant.

\medskip
\noindent
\textbf{Scope of the claim.} We do not claim novelty in the Gaussian
HMM, the sequence-to-sequence LSTM, or the break-even criterion, each of
which is established. The contribution is the framework that connects
them---specifically, the demonstration that label-free state inference
can be made trustworthy enough, and temporally coherent enough, to serve
as the input to an economic control decision on real industrial data, a
connection that Section~\ref{sec:related} shows has not previously been
established.


% ---------------------------------------------------------------------
% 3. CONTRIBUTION LIST
%    Place immediately after the positioning statement. Each item is
%    phrased so that a specific experiment in the paper substantiates
%    it; cross-references point to the sections that do so.
% ---------------------------------------------------------------------

\noindent
The specific contributions of this paper are as follows.

\begin{itemize}
  \item \textbf{A label-free, temporally coherent state inference
  procedure for industrial electrical loads, and an account of what
  temporal coherence actually requires.} Viterbi decoding over a learned
  transition matrix reduces the flicker rate substantially relative to
  i.i.d.\ assignment, but---contrary to the usual assumption---it cannot
  eliminate it, and neither can strengthening the transition prior. The
  penalty a transition matrix can impose on a state change is bounded by
  $\log(p_{\text{self}}/p_{\text{switch}})$, a few nats, whereas
  full-covariance Gaussian emissions fitted to low-noise electrical data
  routinely differ by hundreds of nats between adjacent states, so the
  emission term dominates the decode. We quantify this, show that the
  remedy is an explicit minimum-dwell constraint rather than a stronger
  prior, and demonstrate that the constraint removes the flicker
  essentially without cost to any physical or schedule-based check
  (Section~\ref{sec:method:hmm}, Table~\ref{tab:ablation}).

  \item \textbf{A validation methodology for state partitions in the
  absence of ground truth.} We formulate five mutually independent
  validation tests---power-factor separation, no-load current ratio,
  circuit-level power conservation, dwell-time plausibility, and
  cross-machine replication---none of which uses the model's own
  internal statistics, and aggregate them by consensus
  \cite{strehl2002cluster}. This addresses a problem that recurs
  throughout unsupervised industrial analytics and is, to our knowledge,
  not otherwise treated systematically
  (Section~\ref{sec:validation}).

  \item \textbf{Forecasting of non-productive interval duration from
  inferred states.} We formulate standby-duration prediction as a
  sequence-to-sequence task over label-free inferred states and
  benchmark it against recurrent, convolutional, transformer, and
  tree-based alternatives under a common protocol with statistical
  significance testing (Section~\ref{sec:forecasting}).

  % --- REVISED IN PHASE III -------------------------------------------
  % Previously: "...report break-even sensitivity across electricity
  % tariff and restart energy, giving operators a decision surface."
  % Both halves needed correcting. (i) Break-even is INVARIANT to tariff
  % when the restart cost is purely electrical -- the tariff cancels in
  % Eq. (breakeven) -- so a tariff x restart-energy surface is degenerate;
  % price enters only through non-energy terms. (ii) The restart energy
  % of the studied machine is not measurably positive, so the sensitivity
  % that matters is over the production-delay valuation.
  % --------------------------------------------------------------------
  \item \textbf{An optimisation-based decision layer in which the
  break-even duration is derived rather than assumed, and a
  characterisation of when forecasting is worth its cost.} We replace the
  fixed-threshold shutdown rule with a constrained cost minimisation,
  show that the reduced problem admits an exact offline optimum, and
  measure the restart coefficients from the record instead of assuming
  them---finding that the studied machine has no measurable electrical
  restart penalty, so its restart cost is a production-delay cost. We
  show analytically and numerically that the break-even duration is
  \emph{independent of electricity tariff} whenever the restart cost is
  purely electrical, and therefore report the decision surface over the
  axes that do carry information. Because the idle duration is unknown
  when the decision must be made, the problem is an instance of
  rent-or-buy \cite{karlin1994competitive}; we benchmark against the
  $2$-competitive forecast-free rule rather than against trivial
  policies, and identify the regime---restart cost large relative to the
  idle-duration distribution---in which a forecast earns its place
  (Section~\ref{sec:decision}).

  \item \textbf{End-to-end evaluation on real industrial data, with
  ablation.} We evaluate every stage on IMDELD across multiple machine
  types, quantify each component's marginal contribution by ablation,
  and report recoverable energy, avoided CO\textsubscript{2}, and
  payback period alongside the accuracy metrics
  (Sections~\ref{sec:results}--\ref{sec:sustainability}).

  \item \textbf{Reproducibility.} The full pipeline---preprocessing,
  inference, validation harness, forecasting, and decision layer---is
  released as open source, and all results are obtained from a public
  benchmark dataset.
\end{itemize}


% ---------------------------------------------------------------------
% 4. PAPER ORGANISATION
% ---------------------------------------------------------------------

\noindent
The remainder of this paper is organised as follows.
Section~\ref{sec:related} reviews related work across machine state
identification, non-intrusive load monitoring, industrial energy
optimisation, and predictive shutdown, and states the resulting research
gap. Section~\ref{sec:method} presents the proposed framework.
Section~\ref{sec:validation} describes the ground-truth-free validation
methodology. Section~\ref{sec:results} reports experimental results,
baseline comparisons, and ablation studies.
Section~\ref{sec:sustainability} quantifies the energy, emissions, and
economic implications. Section~\ref{sec:discussion} discusses failure
cases, limitations, and deployment considerations, and
Section~\ref{sec:conclusion} concludes.
baseline comparisons, and ablation studies.
Section~\ref{sec:sustainability} quantifies the energy, emissions, and
economic implications. Section~\ref{sec:discussion} discusses failure
cases, limitations, and deployment considerations, and
Section~\ref{sec:conclusion} concludes.
 after \section{Introduction}.
%  Bibliography: paper/references.bib
% =====================================================================


% ---------------------------------------------------------------------
% 1. REFRAMED ABSTRACT
% ---------------------------------------------------------------------
% Replaces the current abstract. Structure: context -> gap -> what we
% propose (three named components) -> how it is validated -> headline
% result -> significance. Numeric placeholders marked XX must be filled
% from the final experimental run before submission.
%
% PHASE III WARNING ON THE FINAL SENTENCE OF RESULTS.
% The placeholder "...at a payback period of XX years" cannot be filled
% with a favourable figure for pelletizer-I. The plant already
% de-energises for the long idle gaps, leaving 86.8 h of energised idle
% over 63 days of coverage (~325 kWh, ~39 USD/yr at 0.12 USD/kWh), so a
% 5,000 USD retrofit pays back in ~130 years and a three-year payback
% needs a per-machine cost under ~116 USD. Do NOT quote a payback number
% as a headline. The defensible claims are (a) recoverable energy, and
% (b) decision quality: the proposed policy captures 55-76% of the
% attainable head-room and removes the status quo's 20 constraint-
% violating and 29 loss-making shutdowns. Payback belongs in the
% sensitivity discussion. Revisit after Phase IV establishes whether any
% IMDELD machine has a larger energised-idle fraction.
% ---------------------------------------------------------------------

\begin{abstract}
Industrial machines consume a substantial fraction of their total
electrical energy while energised but not producing, yet eliminating
this standby consumption in an existing facility is obstructed by three
coupled difficulties: per-timestamp operating-state annotations do not
exist and are prohibitively expensive to obtain; the duration of a
future non-productive interval is unknown at the moment a shutdown
decision must be made; and shutting a machine down is only economical
when that duration exceeds a break-even point set by restart energy,
tariff, and production constraints. Existing work addresses these in
isolation---non-intrusive load monitoring infers states but does not
forecast or act, while energy-aware scheduling decides optimally but
assumes the idle duration is given. This paper proposes a unified,
label-free framework that closes the loop from raw electrical
measurements to an economically justified shutdown decision. The
framework comprises three components: (i)~\emph{physics-validated
unsupervised state inference}, in which a Gaussian hidden Markov model
jointly learns the emission distributions and the transition dynamics of
the machine's operating states, and the resulting partition---which
cannot be checked against ground truth---is validated by an ensemble of
independent physical, circuit-level, and operational-schedule tests
rather than by supervised metrics; (ii)~\emph{non-productive duration
forecasting}, in which a sequence-to-sequence LSTM predicts the length
of the upcoming standby interval from the inferred state history; and
(iii)~\emph{optimisation-based shutdown decision support}, in which the
break-even duration is derived as the solution of a constrained
cost-minimisation problem over standby, restart, and production-delay
costs, rather than supplied as a tuning parameter. The framework is
evaluated end-to-end on IMDELD, a public benchmark of 1\,Hz
measurements from eight heavy machines in a full-scale industrial
plant. Temporal modelling reduces the state-sequence flicker rate from
XX\% to XX\%, the forecasting stage attains a standby-duration MAE of
XX, and the decision layer identifies XX\,kWh/year of recoverable
standby energy per machine at a payback period of XX~years. The
framework requires only the voltage and current measurements a standard
retrofit meter already provides, and its outputs map directly onto
ISO~50001 energy performance indicators.
\end{abstract}


% ---------------------------------------------------------------------
% 2. POSITIONING / NOVELTY STATEMENT
%    Place at the end of Section 1, immediately after the gap has been
%    stated. This is the sentence reviewers will quote back.
% ---------------------------------------------------------------------

% --- BEFORE (do not use; retained here only to document the change) ---
% "We apply a Gaussian Mixture Model for state identification and an
%  LSTM for forecasting."
%  -> Reads as an application of two off-the-shelf models. Invites the
%     reviewer question "what is new here?" with no defensible answer.
% ----------------------------------------------------------------------

\medskip
\noindent
To close this gap we propose a \emph{unified data-driven framework for
industrial standby energy elimination}. Rather than treating state
inference, forecasting, and control as separate problems, the framework
treats them as one pipeline whose stages are designed against each
other's requirements: the inference stage is constrained to produce a
temporally coherent sequence because duration forecasting is otherwise
ill-posed, and the forecasting stage is constrained to predict duration
rather than class because the decision stage consumes a duration. The
framework consists of three components:

\begin{enumerate}
  \item \textbf{Physics-validated unsupervised state inference.} A
  Gaussian hidden Markov model \cite{rabiner1989tutorial} with a sticky
  self-transition prior \cite{fox2011sticky} jointly estimates the
  emission distributions and the transition dynamics of the machine's
  operating states directly from electrical measurements, requiring no
  annotation. Because no ground-truth labels exist against which the
  resulting partition could be scored, we replace supervised evaluation
  with a validation ensemble drawn from sources external to the model:
  induction-motor power-factor and no-load current physics, conservation
  of power across the plant's circuit hierarchy, dwell-time
  plausibility, cross-machine replication between structurally identical
  units, and agreement with the factory's known shift schedule.

  \item \textbf{Non-productive duration forecasting.} A
  sequence-to-sequence LSTM \cite{hochreiter1997lstm,sutskever2014seq2seq}
  maps the recent inferred-state history to the duration of the upcoming
  standby interval. Predicting duration, rather than the next state
  label, is what makes the output actionable: the break-even criterion
  is a statement about how long a machine will remain idle, not about
  what it is doing now.

  \item \textbf{Optimisation-based shutdown decision support.} The
  shutdown decision is posed as minimisation of total cost---standby
  energy, restart energy and labour, and production-delay penalty---
  subject to thermal cool-down, restart-frequency, and schedule
  constraints, in the spirit of \cite{mouzon2007operational} but with
  the idle duration supplied by component~(ii) instead of by a known
  schedule. The break-even duration is therefore a \emph{derived}
  quantity of the optimisation rather than an operator-supplied
  threshold, and its sensitivity to tariff and restart energy is
  reported explicitly.
\end{enumerate}

\noindent
The framework is validated end-to-end on the IMDELD benchmark
\cite{imdeld2018dataset,martins2018imdeld}, comprising 1\,Hz electrical
measurements from eight heavy machines in a full-scale Brazilian
poultry-feed pellet plant.

\medskip
\noindent
\textbf{Scope of the claim.} We do not claim novelty in the Gaussian
HMM, the sequence-to-sequence LSTM, or the break-even criterion, each of
which is established. The contribution is the framework that connects
them---specifically, the demonstration that label-free state inference
can be made trustworthy enough, and temporally coherent enough, to serve
as the input to an economic control decision on real industrial data, a
connection that Section~\ref{sec:related} shows has not previously been
established.


% ---------------------------------------------------------------------
% 3. CONTRIBUTION LIST
%    Place immediately after the positioning statement. Each item is
%    phrased so that a specific experiment in the paper substantiates
%    it; cross-references point to the sections that do so.
% ---------------------------------------------------------------------

\noindent
The specific contributions of this paper are as follows.

\begin{itemize}
  \item \textbf{A label-free, temporally coherent state inference
  procedure for industrial electrical loads, and an account of what
  temporal coherence actually requires.} Viterbi decoding over a learned
  transition matrix reduces the flicker rate substantially relative to
  i.i.d.\ assignment, but---contrary to the usual assumption---it cannot
  eliminate it, and neither can strengthening the transition prior. The
  penalty a transition matrix can impose on a state change is bounded by
  $\log(p_{\text{self}}/p_{\text{switch}})$, a few nats, whereas
  full-covariance Gaussian emissions fitted to low-noise electrical data
  routinely differ by hundreds of nats between adjacent states, so the
  emission term dominates the decode. We quantify this, show that the
  remedy is an explicit minimum-dwell constraint rather than a stronger
  prior, and demonstrate that the constraint removes the flicker
  essentially without cost to any physical or schedule-based check
  (Section~\ref{sec:method:hmm}, Table~\ref{tab:ablation}).

  \item \textbf{A validation methodology for state partitions in the
  absence of ground truth.} We formulate five mutually independent
  validation tests---power-factor separation, no-load current ratio,
  circuit-level power conservation, dwell-time plausibility, and
  cross-machine replication---none of which uses the model's own
  internal statistics, and aggregate them by consensus
  \cite{strehl2002cluster}. This addresses a problem that recurs
  throughout unsupervised industrial analytics and is, to our knowledge,
  not otherwise treated systematically
  (Section~\ref{sec:validation}).

  \item \textbf{Forecasting of non-productive interval duration from
  inferred states.} We formulate standby-duration prediction as a
  sequence-to-sequence task over label-free inferred states and
  benchmark it against recurrent, convolutional, transformer, and
  tree-based alternatives under a common protocol with statistical
  significance testing (Section~\ref{sec:forecasting}).

  % --- REVISED IN PHASE III -------------------------------------------
  % Previously: "...report break-even sensitivity across electricity
  % tariff and restart energy, giving operators a decision surface."
  % Both halves needed correcting. (i) Break-even is INVARIANT to tariff
  % when the restart cost is purely electrical -- the tariff cancels in
  % Eq. (breakeven) -- so a tariff x restart-energy surface is degenerate;
  % price enters only through non-energy terms. (ii) The restart energy
  % of the studied machine is not measurably positive, so the sensitivity
  % that matters is over the production-delay valuation.
  % --------------------------------------------------------------------
  \item \textbf{An optimisation-based decision layer in which the
  break-even duration is derived rather than assumed, and a
  characterisation of when forecasting is worth its cost.} We replace the
  fixed-threshold shutdown rule with a constrained cost minimisation,
  show that the reduced problem admits an exact offline optimum, and
  measure the restart coefficients from the record instead of assuming
  them---finding that the studied machine has no measurable electrical
  restart penalty, so its restart cost is a production-delay cost. We
  show analytically and numerically that the break-even duration is
  \emph{independent of electricity tariff} whenever the restart cost is
  purely electrical, and therefore report the decision surface over the
  axes that do carry information. Because the idle duration is unknown
  when the decision must be made, the problem is an instance of
  rent-or-buy \cite{karlin1994competitive}; we benchmark against the
  $2$-competitive forecast-free rule rather than against trivial
  policies, and identify the regime---restart cost large relative to the
  idle-duration distribution---in which a forecast earns its place
  (Section~\ref{sec:decision}).

  \item \textbf{End-to-end evaluation on real industrial data, with
  ablation.} We evaluate every stage on IMDELD across multiple machine
  types, quantify each component's marginal contribution by ablation,
  and report recoverable energy, avoided CO\textsubscript{2}, and
  payback period alongside the accuracy metrics
  (Sections~\ref{sec:results}--\ref{sec:sustainability}).

  \item \textbf{Reproducibility.} The full pipeline---preprocessing,
  inference, validation harness, forecasting, and decision layer---is
  released as open source, and all results are obtained from a public
  benchmark dataset.
\end{itemize}


% ---------------------------------------------------------------------
% 4. PAPER ORGANISATION
% ---------------------------------------------------------------------

\noindent
The remainder of this paper is organised as follows.
Section~\ref{sec:related} reviews related work across machine state
identification, non-intrusive load monitoring, industrial energy
optimisation, and predictive shutdown, and states the resulting research
gap. Section~\ref{sec:method} presents the proposed framework.
Section~\ref{sec:validation} describes the ground-truth-free validation
methodology. Section~\ref{sec:results} reports experimental results,
baseline comparisons, and ablation studies.
Section~\ref{sec:sustainability} quantifies the energy, emissions, and
economic implications. Section~\ref{sec:discussion} discusses failure
cases, limitations, and deployment considerations, and
Section~\ref{sec:conclusion} concludes.
baseline comparisons, and ablation studies.
Section~\ref{sec:sustainability} quantifies the energy, emissions, and
economic implications. Section~\ref{sec:discussion} discusses failure
cases, limitations, and deployment considerations, and
Section~\ref{sec:conclusion} concludes.
 after \section{Introduction}.
%  Bibliography: paper/references.bib
% =====================================================================


% ---------------------------------------------------------------------
% 1. REFRAMED ABSTRACT
% ---------------------------------------------------------------------
% Replaces the current abstract. Structure: context -> gap -> what we
% propose (three named components) -> how it is validated -> headline
% result -> significance. Numeric placeholders marked XX must be filled
% from the final experimental run before submission.
%
% PHASE III WARNING ON THE FINAL SENTENCE OF RESULTS.
% The placeholder "...at a payback period of XX years" cannot be filled
% with a favourable figure for pelletizer-I. The plant already
% de-energises for the long idle gaps, leaving 86.8 h of energised idle
% over 63 days of coverage (~325 kWh, ~39 USD/yr at 0.12 USD/kWh), so a
% 5,000 USD retrofit pays back in ~130 years and a three-year payback
% needs a per-machine cost under ~116 USD. Do NOT quote a payback number
% as a headline. The defensible claims are (a) recoverable energy, and
% (b) decision quality: the proposed policy captures 55-76% of the
% attainable head-room and removes the status quo's 20 constraint-
% violating and 29 loss-making shutdowns. Payback belongs in the
% sensitivity discussion. Revisit after Phase IV establishes whether any
% IMDELD machine has a larger energised-idle fraction.
% ---------------------------------------------------------------------

\begin{abstract}
Industrial machines consume a substantial fraction of their total
electrical energy while energised but not producing, yet eliminating
this standby consumption in an existing facility is obstructed by three
coupled difficulties: per-timestamp operating-state annotations do not
exist and are prohibitively expensive to obtain; the duration of a
future non-productive interval is unknown at the moment a shutdown
decision must be made; and shutting a machine down is only economical
when that duration exceeds a break-even point set by restart energy,
tariff, and production constraints. Existing work addresses these in
isolation---non-intrusive load monitoring infers states but does not
forecast or act, while energy-aware scheduling decides optimally but
assumes the idle duration is given. This paper proposes a unified,
label-free framework that closes the loop from raw electrical
measurements to an economically justified shutdown decision. The
framework comprises three components: (i)~\emph{physics-validated
unsupervised state inference}, in which a Gaussian hidden Markov model
jointly learns the emission distributions and the transition dynamics of
the machine's operating states, and the resulting partition---which
cannot be checked against ground truth---is validated by an ensemble of
independent physical, circuit-level, and operational-schedule tests
rather than by supervised metrics; (ii)~\emph{non-productive duration
forecasting}, in which a sequence-to-sequence LSTM predicts the length
of the upcoming standby interval from the inferred state history; and
(iii)~\emph{optimisation-based shutdown decision support}, in which the
break-even duration is derived as the solution of a constrained
cost-minimisation problem over standby, restart, and production-delay
costs, rather than supplied as a tuning parameter. The framework is
evaluated end-to-end on IMDELD, a public benchmark of 1\,Hz
measurements from eight heavy machines in a full-scale industrial
plant. Temporal modelling reduces the state-sequence flicker rate from
XX\% to XX\%, the forecasting stage attains a standby-duration MAE of
XX, and the decision layer identifies XX\,kWh/year of recoverable
standby energy per machine at a payback period of XX~years. The
framework requires only the voltage and current measurements a standard
retrofit meter already provides, and its outputs map directly onto
ISO~50001 energy performance indicators.
\end{abstract}


% ---------------------------------------------------------------------
% 2. POSITIONING / NOVELTY STATEMENT
%    Place at the end of Section 1, immediately after the gap has been
%    stated. This is the sentence reviewers will quote back.
% ---------------------------------------------------------------------

% --- BEFORE (do not use; retained here only to document the change) ---
% "We apply a Gaussian Mixture Model for state identification and an
%  LSTM for forecasting."
%  -> Reads as an application of two off-the-shelf models. Invites the
%     reviewer question "what is new here?" with no defensible answer.
% ----------------------------------------------------------------------

\medskip
\noindent
To close this gap we propose a \emph{unified data-driven framework for
industrial standby energy elimination}. Rather than treating state
inference, forecasting, and control as separate problems, the framework
treats them as one pipeline whose stages are designed against each
other's requirements: the inference stage is constrained to produce a
temporally coherent sequence because duration forecasting is otherwise
ill-posed, and the forecasting stage is constrained to predict duration
rather than class because the decision stage consumes a duration. The
framework consists of three components:

\begin{enumerate}
  \item \textbf{Physics-validated unsupervised state inference.} A
  Gaussian hidden Markov model \cite{rabiner1989tutorial} with a sticky
  self-transition prior \cite{fox2011sticky} jointly estimates the
  emission distributions and the transition dynamics of the machine's
  operating states directly from electrical measurements, requiring no
  annotation. Because no ground-truth labels exist against which the
  resulting partition could be scored, we replace supervised evaluation
  with a validation ensemble drawn from sources external to the model:
  induction-motor power-factor and no-load current physics, conservation
  of power across the plant's circuit hierarchy, dwell-time
  plausibility, cross-machine replication between structurally identical
  units, and agreement with the factory's known shift schedule.

  \item \textbf{Non-productive duration forecasting.} A
  sequence-to-sequence LSTM \cite{hochreiter1997lstm,sutskever2014seq2seq}
  maps the recent inferred-state history to the duration of the upcoming
  standby interval. Predicting duration, rather than the next state
  label, is what makes the output actionable: the break-even criterion
  is a statement about how long a machine will remain idle, not about
  what it is doing now.

  \item \textbf{Optimisation-based shutdown decision support.} The
  shutdown decision is posed as minimisation of total cost---standby
  energy, restart energy and labour, and production-delay penalty---
  subject to thermal cool-down, restart-frequency, and schedule
  constraints, in the spirit of \cite{mouzon2007operational} but with
  the idle duration supplied by component~(ii) instead of by a known
  schedule. The break-even duration is therefore a \emph{derived}
  quantity of the optimisation rather than an operator-supplied
  threshold, and its sensitivity to tariff and restart energy is
  reported explicitly.
\end{enumerate}

\noindent
The framework is validated end-to-end on the IMDELD benchmark
\cite{imdeld2018dataset,martins2018imdeld}, comprising 1\,Hz electrical
measurements from eight heavy machines in a full-scale Brazilian
poultry-feed pellet plant.

\medskip
\noindent
\textbf{Scope of the claim.} We do not claim novelty in the Gaussian
HMM, the sequence-to-sequence LSTM, or the break-even criterion, each of
which is established. The contribution is the framework that connects
them---specifically, the demonstration that label-free state inference
can be made trustworthy enough, and temporally coherent enough, to serve
as the input to an economic control decision on real industrial data, a
connection that Section~\ref{sec:related} shows has not previously been
established.


% ---------------------------------------------------------------------
% 3. CONTRIBUTION LIST
%    Place immediately after the positioning statement. Each item is
%    phrased so that a specific experiment in the paper substantiates
%    it; cross-references point to the sections that do so.
% ---------------------------------------------------------------------

\noindent
The specific contributions of this paper are as follows.

\begin{itemize}
  \item \textbf{A label-free, temporally coherent state inference
  procedure for industrial electrical loads, and an account of what
  temporal coherence actually requires.} Viterbi decoding over a learned
  transition matrix reduces the flicker rate substantially relative to
  i.i.d.\ assignment, but---contrary to the usual assumption---it cannot
  eliminate it, and neither can strengthening the transition prior. The
  penalty a transition matrix can impose on a state change is bounded by
  $\log(p_{\text{self}}/p_{\text{switch}})$, a few nats, whereas
  full-covariance Gaussian emissions fitted to low-noise electrical data
  routinely differ by hundreds of nats between adjacent states, so the
  emission term dominates the decode. We quantify this, show that the
  remedy is an explicit minimum-dwell constraint rather than a stronger
  prior, and demonstrate that the constraint removes the flicker
  essentially without cost to any physical or schedule-based check
  (Section~\ref{sec:method:hmm}, Table~\ref{tab:ablation}).

  \item \textbf{A validation methodology for state partitions in the
  absence of ground truth.} We formulate five mutually independent
  validation tests---power-factor separation, no-load current ratio,
  circuit-level power conservation, dwell-time plausibility, and
  cross-machine replication---none of which uses the model's own
  internal statistics, and aggregate them by consensus
  \cite{strehl2002cluster}. This addresses a problem that recurs
  throughout unsupervised industrial analytics and is, to our knowledge,
  not otherwise treated systematically
  (Section~\ref{sec:validation}).

  \item \textbf{Forecasting of non-productive interval duration from
  inferred states.} We formulate standby-duration prediction as a
  sequence-to-sequence task over label-free inferred states and
  benchmark it against recurrent, convolutional, transformer, and
  tree-based alternatives under a common protocol with statistical
  significance testing (Section~\ref{sec:forecasting}).

  % --- REVISED IN PHASE III -------------------------------------------
  % Previously: "...report break-even sensitivity across electricity
  % tariff and restart energy, giving operators a decision surface."
  % Both halves needed correcting. (i) Break-even is INVARIANT to tariff
  % when the restart cost is purely electrical -- the tariff cancels in
  % Eq. (breakeven) -- so a tariff x restart-energy surface is degenerate;
  % price enters only through non-energy terms. (ii) The restart energy
  % of the studied machine is not measurably positive, so the sensitivity
  % that matters is over the production-delay valuation.
  % --------------------------------------------------------------------
  \item \textbf{An optimisation-based decision layer in which the
  break-even duration is derived rather than assumed, and a
  characterisation of when forecasting is worth its cost.} We replace the
  fixed-threshold shutdown rule with a constrained cost minimisation,
  show that the reduced problem admits an exact offline optimum, and
  measure the restart coefficients from the record instead of assuming
  them---finding that the studied machine has no measurable electrical
  restart penalty, so its restart cost is a production-delay cost. We
  show analytically and numerically that the break-even duration is
  \emph{independent of electricity tariff} whenever the restart cost is
  purely electrical, and therefore report the decision surface over the
  axes that do carry information. Because the idle duration is unknown
  when the decision must be made, the problem is an instance of
  rent-or-buy \cite{karlin1994competitive}; we benchmark against the
  $2$-competitive forecast-free rule rather than against trivial
  policies, and identify the regime---restart cost large relative to the
  idle-duration distribution---in which a forecast earns its place
  (Section~\ref{sec:decision}).

  \item \textbf{End-to-end evaluation on real industrial data, with
  ablation.} We evaluate every stage on IMDELD across multiple machine
  types, quantify each component's marginal contribution by ablation,
  and report recoverable energy, avoided CO\textsubscript{2}, and
  payback period alongside the accuracy metrics
  (Sections~\ref{sec:results}--\ref{sec:sustainability}).

  \item \textbf{Reproducibility.} The full pipeline---preprocessing,
  inference, validation harness, forecasting, and decision layer---is
  released as open source, and all results are obtained from a public
  benchmark dataset.
\end{itemize}


% ---------------------------------------------------------------------
% 4. PAPER ORGANISATION
% ---------------------------------------------------------------------

\noindent
The remainder of this paper is organised as follows.
Section~\ref{sec:related} reviews related work across machine state
identification, non-intrusive load monitoring, industrial energy
optimisation, and predictive shutdown, and states the resulting research
gap. Section~\ref{sec:method} presents the proposed framework.
Section~\ref{sec:validation} describes the ground-truth-free validation
methodology. Section~\ref{sec:results} reports experimental results,
baseline comparisons, and ablation studies.
Section~\ref{sec:sustainability} quantifies the energy, emissions, and
economic implications. Section~\ref{sec:discussion} discusses failure
cases, limitations, and deployment considerations, and
Section~\ref{sec:conclusion} concludes.
baseline comparisons, and ablation studies.
Section~\ref{sec:sustainability} quantifies the energy, emissions, and
economic implications. Section~\ref{sec:discussion} discusses failure
cases, limitations, and deployment considerations, and
Section~\ref{sec:conclusion} concludes.
 after \section{Introduction}.
%  Bibliography: paper/references.bib
% =====================================================================


% ---------------------------------------------------------------------
% 1. REFRAMED ABSTRACT
% ---------------------------------------------------------------------
% Replaces the current abstract. Structure: context -> gap -> what we
% propose (three named components) -> how it is validated -> headline
% result -> significance. Numeric placeholders marked XX must be filled
% from the final experimental run before submission.
%
% PHASE III WARNING ON THE FINAL SENTENCE OF RESULTS.
% The placeholder "...at a payback period of XX years" cannot be filled
% with a favourable figure for pelletizer-I. The plant already
% de-energises for the long idle gaps, leaving 86.8 h of energised idle
% over 63 days of coverage (~325 kWh, ~39 USD/yr at 0.12 USD/kWh), so a
% 5,000 USD retrofit pays back in ~130 years and a three-year payback
% needs a per-machine cost under ~116 USD. Do NOT quote a payback number
% as a headline. The defensible claims are (a) recoverable energy, and
% (b) decision quality: the proposed policy captures 55-76% of the
% attainable head-room and removes the status quo's 20 constraint-
% violating and 29 loss-making shutdowns. Payback belongs in the
% sensitivity discussion. Revisit after Phase IV establishes whether any
% IMDELD machine has a larger energised-idle fraction.
% ---------------------------------------------------------------------

\begin{abstract}
Industrial machines consume a substantial fraction of their total
electrical energy while energised but not producing, yet eliminating
this standby consumption in an existing facility is obstructed by three
coupled difficulties: per-timestamp operating-state annotations do not
exist and are prohibitively expensive to obtain; the duration of a
future non-productive interval is unknown at the moment a shutdown
decision must be made; and shutting a machine down is only economical
when that duration exceeds a break-even point set by restart energy,
tariff, and production constraints. Existing work addresses these in
isolation---non-intrusive load monitoring infers states but does not
forecast or act, while energy-aware scheduling decides optimally but
assumes the idle duration is given. This paper proposes a unified,
label-free framework that closes the loop from raw electrical
measurements to an economically justified shutdown decision. The
framework comprises three components: (i)~\emph{physics-validated
unsupervised state inference}, in which a Gaussian hidden Markov model
jointly learns the emission distributions and the transition dynamics of
the machine's operating states, and the resulting partition---which
cannot be checked against ground truth---is validated by an ensemble of
independent physical, circuit-level, and operational-schedule tests
rather than by supervised metrics; (ii)~\emph{non-productive duration
forecasting}, in which a sequence-to-sequence LSTM predicts the length
of the upcoming standby interval from the inferred state history; and
(iii)~\emph{optimisation-based shutdown decision support}, in which the
break-even duration is derived as the solution of a constrained
cost-minimisation problem over standby, restart, and production-delay
costs, rather than supplied as a tuning parameter. The framework is
evaluated end-to-end on IMDELD, a public benchmark of 1\,Hz
measurements from eight heavy machines in a full-scale industrial
plant. Temporal modelling reduces the state-sequence flicker rate from
XX\% to XX\%, the forecasting stage attains a standby-duration MAE of
XX, and the decision layer identifies XX\,kWh/year of recoverable
standby energy per machine at a payback period of XX~years. The
framework requires only the voltage and current measurements a standard
retrofit meter already provides, and its outputs map directly onto
ISO~50001 energy performance indicators.
\end{abstract}


% ---------------------------------------------------------------------
% 2. POSITIONING / NOVELTY STATEMENT
%    Place at the end of Section 1, immediately after the gap has been
%    stated. This is the sentence reviewers will quote back.
% ---------------------------------------------------------------------

% --- BEFORE (do not use; retained here only to document the change) ---
% "We apply a Gaussian Mixture Model for state identification and an
%  LSTM for forecasting."
%  -> Reads as an application of two off-the-shelf models. Invites the
%     reviewer question "what is new here?" with no defensible answer.
% ----------------------------------------------------------------------

\medskip
\noindent
To close this gap we propose a \emph{unified data-driven framework for
industrial standby energy elimination}. Rather than treating state
inference, forecasting, and control as separate problems, the framework
treats them as one pipeline whose stages are designed against each
other's requirements: the inference stage is constrained to produce a
temporally coherent sequence because duration forecasting is otherwise
ill-posed, and the forecasting stage is constrained to predict duration
rather than class because the decision stage consumes a duration. The
framework consists of three components:

\begin{enumerate}
  \item \textbf{Physics-validated unsupervised state inference.} A
  Gaussian hidden Markov model \cite{rabiner1989tutorial} with a sticky
  self-transition prior \cite{fox2011sticky} jointly estimates the
  emission distributions and the transition dynamics of the machine's
  operating states directly from electrical measurements, requiring no
  annotation. Because no ground-truth labels exist against which the
  resulting partition could be scored, we replace supervised evaluation
  with a validation ensemble drawn from sources external to the model:
  induction-motor power-factor and no-load current physics, conservation
  of power across the plant's circuit hierarchy, dwell-time
  plausibility, cross-machine replication between structurally identical
  units, and agreement with the factory's known shift schedule.

  \item \textbf{Non-productive duration forecasting.} A
  sequence-to-sequence LSTM \cite{hochreiter1997lstm,sutskever2014seq2seq}
  maps the recent inferred-state history to the duration of the upcoming
  standby interval. Predicting duration, rather than the next state
  label, is what makes the output actionable: the break-even criterion
  is a statement about how long a machine will remain idle, not about
  what it is doing now.

  \item \textbf{Optimisation-based shutdown decision support.} The
  shutdown decision is posed as minimisation of total cost---standby
  energy, restart energy and labour, and production-delay penalty---
  subject to thermal cool-down, restart-frequency, and schedule
  constraints, in the spirit of \cite{mouzon2007operational} but with
  the idle duration supplied by component~(ii) instead of by a known
  schedule. The break-even duration is therefore a \emph{derived}
  quantity of the optimisation rather than an operator-supplied
  threshold, and its sensitivity to tariff and restart energy is
  reported explicitly.
\end{enumerate}

\noindent
The framework is validated end-to-end on the IMDELD benchmark
\cite{imdeld2018dataset,martins2018imdeld}, comprising 1\,Hz electrical
measurements from eight heavy machines in a full-scale Brazilian
poultry-feed pellet plant.

\medskip
\noindent
\textbf{Scope of the claim.} We do not claim novelty in the Gaussian
HMM, the sequence-to-sequence LSTM, or the break-even criterion, each of
which is established. The contribution is the framework that connects
them---specifically, the demonstration that label-free state inference
can be made trustworthy enough, and temporally coherent enough, to serve
as the input to an economic control decision on real industrial data, a
connection that Section~\ref{sec:related} shows has not previously been
established.


% ---------------------------------------------------------------------
% 3. CONTRIBUTION LIST
%    Place immediately after the positioning statement. Each item is
%    phrased so that a specific experiment in the paper substantiates
%    it; cross-references point to the sections that do so.
% ---------------------------------------------------------------------

\noindent
The specific contributions of this paper are as follows.

\begin{itemize}
  \item \textbf{A label-free, temporally coherent state inference
  procedure for industrial electrical loads, and an account of what
  temporal coherence actually requires.} Viterbi decoding over a learned
  transition matrix reduces the flicker rate substantially relative to
  i.i.d.\ assignment, but---contrary to the usual assumption---it cannot
  eliminate it, and neither can strengthening the transition prior. The
  penalty a transition matrix can impose on a state change is bounded by
  $\log(p_{\text{self}}/p_{\text{switch}})$, a few nats, whereas
  full-covariance Gaussian emissions fitted to low-noise electrical data
  routinely differ by hundreds of nats between adjacent states, so the
  emission term dominates the decode. We quantify this, show that the
  remedy is an explicit minimum-dwell constraint rather than a stronger
  prior, and demonstrate that the constraint removes the flicker
  essentially without cost to any physical or schedule-based check
  (Section~\ref{sec:method:hmm}, Table~\ref{tab:ablation}).

  \item \textbf{A validation methodology for state partitions in the
  absence of ground truth.} We formulate five mutually independent
  validation tests---power-factor separation, no-load current ratio,
  circuit-level power conservation, dwell-time plausibility, and
  cross-machine replication---none of which uses the model's own
  internal statistics, and aggregate them by consensus
  \cite{strehl2002cluster}. This addresses a problem that recurs
  throughout unsupervised industrial analytics and is, to our knowledge,
  not otherwise treated systematically
  (Section~\ref{sec:validation}).

  \item \textbf{Forecasting of non-productive interval duration from
  inferred states.} We formulate standby-duration prediction as a
  sequence-to-sequence task over label-free inferred states and
  benchmark it against recurrent, convolutional, transformer, and
  tree-based alternatives under a common protocol with statistical
  significance testing (Section~\ref{sec:forecasting}).

  % --- REVISED IN PHASE III -------------------------------------------
  % Previously: "...report break-even sensitivity across electricity
  % tariff and restart energy, giving operators a decision surface."
  % Both halves needed correcting. (i) Break-even is INVARIANT to tariff
  % when the restart cost is purely electrical -- the tariff cancels in
  % Eq. (breakeven) -- so a tariff x restart-energy surface is degenerate;
  % price enters only through non-energy terms. (ii) The restart energy
  % of the studied machine is not measurably positive, so the sensitivity
  % that matters is over the production-delay valuation.
  % --------------------------------------------------------------------
  \item \textbf{An optimisation-based decision layer in which the
  break-even duration is derived rather than assumed, and a
  characterisation of when forecasting is worth its cost.} We replace the
  fixed-threshold shutdown rule with a constrained cost minimisation,
  show that the reduced problem admits an exact offline optimum, and
  measure the restart coefficients from the record instead of assuming
  them---finding that the studied machine has no measurable electrical
  restart penalty, so its restart cost is a production-delay cost. We
  show analytically and numerically that the break-even duration is
  \emph{independent of electricity tariff} whenever the restart cost is
  purely electrical, and therefore report the decision surface over the
  axes that do carry information. Because the idle duration is unknown
  when the decision must be made, the problem is an instance of
  rent-or-buy \cite{karlin1994competitive}; we benchmark against the
  $2$-competitive forecast-free rule rather than against trivial
  policies, and identify the regime---restart cost large relative to the
  idle-duration distribution---in which a forecast earns its place
  (Section~\ref{sec:decision}).

  \item \textbf{End-to-end evaluation on real industrial data, with
  ablation.} We evaluate every stage on IMDELD across multiple machine
  types, quantify each component's marginal contribution by ablation,
  and report recoverable energy, avoided CO\textsubscript{2}, and
  payback period alongside the accuracy metrics
  (Sections~\ref{sec:results}--\ref{sec:sustainability}).

  \item \textbf{Reproducibility.} The full pipeline---preprocessing,
  inference, validation harness, forecasting, and decision layer---is
  released as open source, and all results are obtained from a public
  benchmark dataset.
\end{itemize}


% ---------------------------------------------------------------------
% 4. PAPER ORGANISATION
% ---------------------------------------------------------------------

\noindent
The remainder of this paper is organised as follows.
Section~\ref{sec:related} reviews related work across machine state
identification, non-intrusive load monitoring, industrial energy
optimisation, and predictive shutdown, and states the resulting research
gap. Section~\ref{sec:method} presents the proposed framework.
Section~\ref{sec:validation} describes the ground-truth-free validation
methodology. Section~\ref{sec:results} reports experimental results,
baseline comparisons, and ablation studies.
Section~\ref{sec:sustainability} quantifies the energy, emissions, and
economic implications. Section~\ref{sec:discussion} discusses failure
cases, limitations, and deployment considerations, and
Section~\ref{sec:conclusion} concludes.
baseline comparisons, and ablation studies.
Section~\ref{sec:sustainability} quantifies the energy, emissions, and
economic implications. Section~\ref{sec:discussion} discusses failure
cases, limitations, and deployment considerations, and
Section~\ref{sec:conclusion} concludes.
